\section*{Chapitre \ref{ch:b3:green-industrial} --- Chimie verte et procédés industriels}

\begin{solution}{exo:b3:green-industrial:1}
Diels--Alder~: 100\,\% (une addition). Estérification~: $88.0/(60.0 + 46.0) = 83$\,\%. Wittig~:
$104.0/(106.0 + 276.0) = 27$\,\%, l'oxyde de triphénylphosphine emportant l'essentiel de la masse.
\end{solution}

\begin{solution}{exo:b3:green-industrial:2}
$\mathrm{PMI} = 120/8.0 = 15$~; $E = \mathrm{PMI} - 1 = 14$.
\end{solution}

\begin{solution}{exo:b3:green-industrial:3}
Un lavage d'une charge normalisée à une température et avec un résultat de
lavage donnés. Les lessives se dosent différemment~: une poudre concentrée
en exige moins par lavage, si bien que des masses égales ne rendent pas le
même service.
\end{solution}

\begin{solution}{exo:b3:green-industrial:4}
Le 2-méthyltétrahydrofurane ou l'acétate d'éthyle à la place du
dichlorométhane~; l'heptane (avec l'acétate d'éthyle) à la place de
l'hexane~; pour le DMF, l'acétate d'éthyle, le 2-méthyltétrahydrofurane ou
le carbonate de diméthyle, ou un couplage dans l'eau avec des \omterm{def:b3:colloids:surfactant}{tensioactifs},
selon les réactifs.
\end{solution}

\begin{solution}{exo:b3:green-industrial:5}
1.000 mol d'acide~; 0.800 mol d'ester (\qty{88.0}{g/mol})~: rendement de 80\,\%. AE 83.0\,\%. SF $= 152/106 = 1.43$.
$\mathrm{RME} = 70.4/152 = 0.463$~; $\mathrm{AE} \times \text{rendement}/\mathrm{SF} = 0.830 \times 0.80/1.43 = 0.463$.
\end{solution}

\begin{solution}{exo:b3:green-industrial:6}
Chlorhydrine~: $44.0/(28.0 + 71.0 + 74.1) = 25$\,\%~; oxydation directe~: 100\,\%. Une mole de $\ce{CaCl2}$ par mole
d'oxyde d'éthylène~: $\qty{1.0e6}{g}/\qty{44.0}{g/mol} \times \qty{111.1}{g/mol} = \qty{2.5}{t}$ par tonne.
\end{solution}

\begin{solution}{exo:b3:green-industrial:7}
$3/8 \times \qty{5.88e4}{mol} \times \qty{44.0}{g/mol} = \qty{0.97}{t}$ de $\ce{CO2}$ par tonne d'ammoniac. Pour le dihydrogène, $\ce{CH4 + 2H2O -> CO2 + 4H2}$~:
$1/4$ mole par mole, $\qty{5.0e5}{mol}/4 \times \qty{44.0}{g/mol} = \qty{5.5}{t}$ de $\ce{CO2}$ par tonne de dihydrogène.
\end{solution}

\begin{solution}{exo:b3:green-industrial:8}
% ledger: dfg:H2O_l, dfh:H2O-l
Un kilogramme représente \qty{500}{mol}. À partir de $\Delta_rG^\circ = \qty{237.1}{kJ/mol}$~: $\qty{1.19e5}{kJ} = \qty{32.9}{kWh}$. À partir de $\Delta_rH^\circ = \qty{285.8}{kJ/mol}$~: \qty{39.7}{kWh}
(la différence est de la chaleur que le milieu extérieur peut fournir).
\end{solution}

\begin{solution}{exo:b3:green-industrial:9}
$\qty{1.0e6}{g}/\qty{146.0}{g/mol} = \qty{6.85e3}{mol}$ de $\ce{N2O}$, $\qty{0.30}{t}$~; $\times 273 = \qty{82}{t}$ d'équivalent $\ce{CO2}$ par tonne d'acide adipique.
\end{solution}

\begin{solution}{exo:b3:green-industrial:10}
La RME d'une étape compte le réactif en excès comme consommé~; si l'excès
est récupéré et réintroduit, seule la part réellement perdue est
consommée, et l'efficacité massique globale est supérieure au produit des
valeurs des étapes. Dans le \omterm{def:b3:green-industrial:e-factor}{facteur E}, le flux recyclé n'est pas un
déchet, à condition que le recyclage soit à l'intérieur du périmètre~; il
faut alors compter son énergie et ses pertes.
\end{solution}

\begin{solution}{exo:b3:green-industrial:11}
Avant~: $20 \times 0.80 = \qty{16}{kg}$ de solvant, dont 10\,\% perdus~: \qty{1.6}{kg} de déchets. Après~: \qty{4.0}{kg}, dont \qty{0.40}{kg} perdus. Le
\omterm{def:b3:green-industrial:e-factor}{facteur E} complet diminue de 1.2 kg par kilogramme de produit (et l'énergie
de distillation de trois quarts).
\end{solution}

\begin{solution}{exo:b3:green-industrial:12}
Par \omterm{def:b3:green-industrial:lca}{unité fonctionnelle} et avec le même périmètre, elle présenterait
séparément chaque catégorie d'impact (climat, usage des sols, consommation
d'eau\dots) plutôt qu'un score unique. Une décision exige les poids
attribués à ces catégories, la rareté locale de l'eau et des terres,
l'incertitude des données, et de savoir si la matière première biosourcée
entre en concurrence avec l'alimentation.
\end{solution}

\begin{solution}{pb:b3:green-industrial:1}
\textbf{1.} $\ce{C10H14 + C4H6O3 -> C12H16O + C2H4O2}$ (HF)~; $\ce{C12H16O + H2 -> C12H18O}$ (nickel de Raney)~;
$\ce{C12H18O + CO -> C13H18O2}$ (palladium).

\textbf{2.} $\ce{C10H14 + C4H6O3 + H2 + CO -> C13H18O2 + C2H4O2}$.

\textbf{3.} $134.0 + 102.0 + 2.0 + 28.0 = \qty{266.0}{g/mol}$ de réactifs~; ibuprofène \qty{206.0}{g/mol}~: AE $= 206.0/266.0 = 77$\,\%.

\textbf{4.} Isobutylbenzène, anhydride acétique, chloroacétate d'éthyle,
éthanolate de sodium, acide aqueux ($\ce{H3O+}$), hydroxylamine et eau.

\textbf{5.} $134.0 + 102.0 + 122.5 + 68.0 + 19.0 + 33.0 + 2 \times 18.0 = \qty{514.5}{g/mol}$~: AE $= 206.0/514.5 = 40$\,\%.

\textbf{6.} Acide acétique, chlorure de sodium, éthanol, dioxyde de carbone,
eau, ammoniac (sous forme de sel d'ammonium), et sels d'aluminium issus du
chlorure d'aluminium stœchiométrique.

\textbf{7.} $0.71 + 0.55 + 0.010 + 0.14 + 0.05 = \qty{1.46}{t}$.

\textbf{8.} PMI 1.46~; $E = 0.46$.

\textbf{9.} $0.46 - 0.31 = 0.15$.

\textbf{10.} PMI $= 2.0 + 1.5 = 3.5$~; $E = 2.5$.

\textbf{11.} Les rendements inférieurs à 100\,\%, les excès de réactifs, les
pertes de solvant et de catalyseur, et les matières du traitement ajoutent
tous des déchets que l'\omterm{def:b3:green-industrial:atom-economy}{économie d'atomes} ignore.

\textbf{12.} Prévenir les déchets, maximiser l'\omterm{def:b3:green-industrial:atom-economy}{économie d'atomes}, utiliser
des catalyseurs plutôt que des réactifs stœchiométriques, éviter les
dérivés et les étapes inutiles, économiser l'énergie.

\textbf{13.} Le chlorure d'aluminium stœchiométrique, hydrolysé en déchets
d'aluminium lors du traitement~; le fluorure d'hydrogène est à la fois
catalyseur et solvant, et il est récupéré et réutilisé.

\textbf{14.} Le nickel de Raney (un catalyseur d'hydrogénation hétérogène).

\textbf{15.} Un complexe phosphine du palladium~: la chimie de \omterm{def:b3:homogeneous-catalysis:carbonylation}{carbonylation} du
\cref{ch:b3:homogeneous-catalysis}.

\textbf{16.} L'anhydride acétique transfère un groupe acétyle~; l'autre
moitié de la molécule part sous forme d'acide acétique.

\textbf{17.} Le fluorure d'hydrogène est très toxique et corrosif, le
monoxyde de carbone toxique et inflammable~: équipements clos en matériaux
résistants, détecteurs et opérateurs formés.

\textbf{18.} Chaque étape ajoute chauffages, refroidissements, séparations
et récupération de solvants~; trois étapes en exigent moins.

\textbf{19.} $2.5 \times 3000 = \qty{7500}{t}$ par an.

\textbf{20.} $0.15 \times 3000 = \qty{450}{t}$ par an.

\textbf{21.} Environ \qty{7000}{t} de déchets évités chaque année.

\textbf{22.} Non~: la nature des déchets compte (une tonne d'eau salée n'est
pas une tonne d'un composé toxique persistant), de même que l'énergie, les
émissions et les dangers des réactifs.

\textbf{23.} Les impacts de la fabrication des réactifs et de l'énergie en
amont, des émissions dans l'air et dans l'eau, et de l'utilisation et de
l'élimination du produit, par \omterm{def:b3:green-industrial:lca}{unité fonctionnelle}, dans plusieurs
catégories d'impact.

\textbf{24.} La voie catalytique en trois étapes a une \omterm{def:b3:green-industrial:atom-economy}{économie d'atomes} d'environ 77\,\% (40\,\% pour la
voie en six étapes).
\end{solution}
